\section{AS-Depth Design}
\label{sec:design}

\system{} combines three mechanisms derived from the opportunities in
Section~\ref{sec:motivation}.  A task-level execution framework partitions the
model and overlaps consecutive frames; DAQ coordinates low precision across
the Thor encoder and FPGA decoder; and a dedicated decoder accelerator turns
the regular decoder workload into an efficient hardware dataflow.

\begin{figure}[t]
  \centering
  \figslot{figures/as-depth-system.pdf}{1.55in}{%
    AS-Depth system architecture\\[5pt]
    Thor: RGB-D preprocessing $\rightarrow$ encoder $\rightarrow$ quantized boundary feature\\
    Queues / double buffers / Ethernet $\leftrightarrow$ FPGA decoder\\
    Inset: cross-frame pipeline for frames $i$, $i+1$, and $i+2$}
  \caption{Overview of \system.  Thor executes the representation encoder and
  produces the cross-device feature tensor; the FPGA executes the depth
  decoder.  Queues and double buffers preserve frame order while overlapping
  encoding, communication, and decoding across consecutive frames.}
  \Description{A system block diagram of the Thor encoder, quantized feature
  boundary, communication queues, FPGA decoder, completed-depth output, and a
  cross-frame pipeline inset.}
  \label{fig:system-overview}
\end{figure}

\begin{figure}[t]
  \centering
  \figslot{figures/fpga-decoder-accelerator.pdf}{1.70in}{%
    Dedicated FPGA depth-decoder accelerator\\[5pt]
    Ethernet/interface engine $\rightarrow$ command scheduler $\rightarrow$ compute engines\\
    Weight/activation buffers $\leftrightarrow$ external memory\\
    Label weighted, residual, activation, and upsampling paths}
  \caption{Architecture of the dedicated FPGA depth-decoder accelerator.  The
  final figure will show the communication front end, programmable layer
  controller, weighted and auxiliary compute engines, local reuse buffers, and
  external-memory datapath.}
  \Description{A block diagram of the FPGA decoder accelerator and its
  communication, scheduling, compute, buffering, and memory components.}
  \label{fig:fpga-accelerator}
\end{figure}

\subsection{Task-Level Asynchronous Heterogeneous Execution Framework}
\label{sec:framework}

\paragraph{Model partition and tensor contract.}
The model is partitioned once at the encoder--decoder boundary.  Thor consumes
the RGB image and sensor depth, executes the encoder, and exports a quantized
feature tensor.  The FPGA consumes this tensor and returns the completed
metric-depth map.  The final paper will state the exact tensor shape, layout,
precision, scale metadata, and frame identifier used by this interface.

\paragraph{Queued cross-frame execution.}
The host runtime assigns separate buffers to encoder output, outbound
transfer, decoder input, and completed depth.  Producer--consumer queues make
ownership explicit and prevent a buffer from being reused before its consumer
finishes.  After pipeline fill, the encoder for the next frame can progress in
parallel with transfer and decoder execution for the current frame.  We report
both single-frame latency and steady-state throughput because only the latter
benefits fully from cross-frame overlap.

\paragraph{Communication and back pressure.}
The runtime batches feature metadata with the boundary tensor and preserves
frame order on the return path.  If either device becomes the bottleneck,
bounded queues apply back pressure instead of silently dropping or reordering
frames.  Section~\ref{sec:evaluation} measures transfer time separately from
useful computation so that communication hiding can be verified.

\subsection{Depth-Aware Quantization and Customized Precision}
\label{sec:daq}

DAQ treats quantization as an end-to-end property of the heterogeneous model.
On Thor, it targets selected MLP computation in the encoder; at the partition
boundary, it defines the representation transmitted to the FPGA; and inside
the FPGA, it determines the precision of decoder weights, activations,
accumulators, and output depth.  Training optimizes the quantized path against
metric-depth supervision so that numerical error is weighted by its effect on
the final depth map rather than by feature reconstruction alone, extending the
standard quantization-aware training setting~\cite{jacob2018quantization}.

\draftnote{Specify the encoder MLP layers covered by DAQ, the exact bit widths,
per-tensor or per-channel scaling, the depth-aware loss, calibration procedure,
and any layers retained at higher precision.  The ``without DAQ'' ablation must
keep the same partition and target precision while replacing DAQ with a
conventional quantization policy.}

\begin{figure}[t]
  \centering
  \figslot{figures/daq-precision.pdf}{1.42in}{%
    DAQ and precision map\\[5pt]
    Encoder MLP $\rightarrow$ boundary tensor $\rightarrow$ decoder\\
    Mark bit width, scale, accumulator, and depth-aware loss}
  \caption{Depth-aware quantization across the heterogeneous inference path.
  The final diagram will identify the quantized encoder MLP operations, the
  cross-device representation, decoder precision, and the depth-domain
  training signal.}
  \Description{A precision-flow diagram linking quantized encoder MLP layers,
  the transferred feature tensor, the FPGA decoder, and a depth-aware loss.}
  \label{fig:daq}
\end{figure}

\subsection{Dedicated FPGA Depth Decoder Accelerator}
\label{sec:accelerator}

The FPGA accelerator specializes the weighted and auxiliary operators in the
depth decoder while retaining a programmable layer schedule.  Its datapath
must expose the precision contract established by DAQ, reuse weights and
activations on chip where profitable, and stream larger intermediate tensors
through external memory without changing model semantics.  The communication
front end converts received feature packets into accelerator-visible buffers
and returns the completed depth map after the final layer.  This design follows
the broader principle of mapping quantized networks to customized FPGA
dataflows~\cite{umuroglu2017finn} while specializing the execution path for
metric-depth decoding.

\paragraph{Compute and dataflow.}
The architecture description will quantify parallel multiply--accumulate
lanes, accumulator precision, supported convolution variants, and the mapping
of activation, residual, and upsampling operators.  A layer controller selects
the appropriate engine and coordinates local buffers with external memory.

\paragraph{Control and correctness boundary.}
For every layer, the scheduler records tensor dimensions, addresses,
quantization parameters, and dependencies.  Verification compares each
hardware layer boundary and the final depth map against the quantized software
reference before end-to-end performance is reported.
